\chapter{Data Preparation and Feature Engineering}
\label{ch:dataprep}
In practice, more model quality comes from features and clean data than from the choice of algorithm, and
most production bugs live here. OA papers test the crisp parts (which models need scaling, target versus
one-hot encoding, IDF arithmetic, missing values); interviews test judgement (``this feature is 40\% missing,
what do you do?'').

\section{Scaling}
\prototype{Wine data: $k$-NN at 66\% accuracy before standardising, 96\% after.}
\begin{itemize}
	\ii \vocab{Standardisation} $(x-\mu)/\sigma$: zero mean, unit variance; the default.
	\ii \vocab{Min--max} $(x-\min)/(\max-\min)$ to $[0,1]$: bounded inputs (pixels, some neural nets); one
	outlier squashes everything else.
	\ii \vocab{Robust} $(x-\text{median})/\text{IQR}$: resistant to outliers.
\end{itemize}
For the column $10,12,11,13,12,100$:
\begin{itemize}
	\ii standardised: $\gv{c15_dataprep/sc_std}$;
	\ii min--max: $\gv{c15_dataprep/sc_mm}$ --- the five ordinary values are crushed near 0;
	\ii robust: $\gv{c15_dataprep/sc_rob}$ --- they stay spread and the outlier is obviously extreme.
\end{itemize} Log or Box--Cox transforms handle skew before
scaling.

\begin{example}[Who needs scaling]
	5-fold accuracy on the wine data, raw versus standardised inside a pipeline:
	\begin{center}
		\small
		\begin{tabular}{lcc}
			\toprule
			model & raw & standardised \\ \midrule
			\gv{c15_dataprep/scale_rows}
			\bottomrule
		\end{tabular}
	\end{center}
	Distance- and margin-based models ($k$-NN, SVM, $k$-means), variance-based PCA, and penalised or
	gradient-trained linear models depend on scale (logistic regression also needed far more iterations
	unscaled). Trees and tree ensembles split on thresholds and are invariant to any monotone rescaling.
	Unpenalised OLS gives the same predictions either way.
\end{example}

\section{Encoding categoricals}
\prototype{A ``city'' column with 1500 distinct values and a churn label.}
\begin{itemize}
	\ii \vocab{One-hot}: one binary column per level (drop one to avoid the dummy-variable trap in unpenalised
	linear models). Safe and interpretable; dimension explodes with cardinality and rare levels are noisy.
	\ii \vocab{Ordinal/label encoding}: integers $0,1,2,\dots$. Fine for trees and for genuinely ordered
	categories (small $<$ medium $<$ large); for linear models it invents a spurious order and spacing.
	\ii \vocab{Frequency encoding}, \vocab{feature hashing} (hash levels into a fixed number of columns;
	collisions are the price of bounded memory, useful for huge or streaming vocabularies), learned
	\vocab{embeddings} (neural nets).
	\ii \vocab{Target (mean) encoding}: replace a level by a (smoothed) mean of the target over rows with that
	level --- one dense column regardless of cardinality.
\end{itemize}

\subsection{Target encoding and its leakage}
With $n_c$ rows in level $c$, level mean $\bar y_c$, global mean $\bar y$ and smoothing $m$:
\[ \mathrm{TE}(c)=\frac{n_c\bar y_c+m\,\bar y}{n_c+m}, \]
a shrinkage towards the global mean that is strong for rare levels (the Beta-prior estimate of
\Cref{ch:maths} again). A level seen 3 times, always positive, with $\bar y=\gv{c15_dataprep/te_gmean}$ and
$m=\gv{c15_dataprep/te_m}$, encodes to $\gv{c15_dataprep/te_a}$, not $1$ (matches scikit-learn's
\code{TargetEncoder(smooth=10)}).

\begin{example}[The leak, measured]
	4000 rows; the label depends only on a numeric feature $x$; a 1500-level ``ID-like'' category is pure noise.
	Logistic regression on $x$ alone: test AUC $\gv{c15_dataprep/base_te_auc}$. Add the category, target-encoded
	with means computed on the full training set (each row's own label included): \emph{training} AUC
	$\gv{c15_dataprep/leak_tr_auc}$, coefficient on the encoding $\gv{c15_dataprep/leak_coef}$, \emph{test} AUC
	$\gv{c15_dataprep/leak_te_auc}$ --- worse than ignoring the column. With levels of about three rows, the
	encoding largely \emph{is} the row's own label. Cross-fitted (out-of-fold) encoding, as
	\code{TargetEncoder.fit_transform} does, gives training AUC $\gv{c15_dataprep/oof_tr_auc}$, test AUC
	$\gv{c15_dataprep/oof_te_auc}$ and a near-zero coefficient ($\gv{c15_dataprep/oof_coef}$): the noise is
	correctly ignored.
\end{example}

\begin{remark}
	\trap ``Target encoding always beats one-hot for high cardinality.'' Only if done out-of-fold (or with
	CatBoost's ordered statistics) and with smoothing; naive target encoding is a textbook leak that inflates
	validation scores and hurts production. One-hot plus regularisation is the safe baseline; for trees, ordinal
	or native categorical handling is often enough.
\end{remark}

\section{Missing values and outliers}
\prototype{28\% of applicants skip the income question --- mostly high earners.}
Missingness mechanisms: \vocab{MCAR} (completely at random: deleting rows only costs data), \vocab{MAR}
(depends on observed features: model-based imputation works), \vocab{MNAR} (depends on the missing value itself:
the fact of missingness is information). Options: drop rows (small, MCAR), drop columns (mostly empty and
uninformative), mean/median/mode imputation (fast, shrinks variance and distorts correlations), $k$-NN or
iterative model-based imputation, a \emph{missing-indicator} column, or models that handle \texttt{NaN} natively
(XGBoost, LightGBM, recent scikit-learn trees). Fit imputers on training folds only.

\begin{example}[Missingness is a feature]
	Income missing for $\gv{c15_dataprep/miss_frac}$ of applicants, more often for high earners; the target rises
	with income. Mean imputation alone: test AUC $\gv{c15_dataprep/miss_auc_mean}$ (and the income standard
	deviation drops from $\gv{c15_dataprep/miss_sd_before}$ to $\gv{c15_dataprep/miss_sd_after}$). Mean imputation
	plus an \code{is_missing} indicator: $\gv{c15_dataprep/miss_auc_ind}$.
\end{example}

\vocab{Outliers}: flag with the IQR rule (outside $[Q_1-1.5\,\mathrm{IQR},\,Q_3+1.5\,\mathrm{IQR}]$; for the column
above $Q_1=\gv{c15_dataprep/iqr_q1}$, $Q_3=\gv{c15_dataprep/iqr_q3}$, upper fence $\gv{c15_dataprep/iqr_hi}$, so 100 is
flagged), $z$-scores (fragile: the outlier inflates $\sigma$), or isolation forests (\Cref{ch:unsup}). Then decide:
data error (fix or drop), genuine extreme (keep; use robust losses, log transforms or winsorising), or the
signal itself (fraud).

\section{Feature selection}
\prototype{50 features, 5 informative: which methods find them?}
\vocab{Filter} methods score features one at a time (correlation, ANOVA $F$, mutual information, $\chi^2$): fast,
ignore interactions and redundancy. \vocab{Wrapper} methods search subsets using a model (forward/backward
selection, recursive feature elimination): expensive, can overfit the validation data. \vocab{Embedded} methods
select while fitting (lasso, tree importances). On 600 rows with 5 informative of 50 features: mutual
information picks $\gv{c15_dataprep/fs_mi_hit}$ of the 5 in its top 5, RFE with logistic regression
$\gv{c15_dataprep/fs_rfe_hit}$ of 5, and an $\ell_1$ logistic regression keeps $\gv{c15_dataprep/fs_l1_n}$ features
including $\gv{c15_dataprep/fs_l1_hit}$ true ones. Any selection that looks at labels must run inside the CV loop
(\Cref{ch:eval}).

\section{Imbalance}
\prototype{5\% positives; SMOTE before cross-validation reports F1 0.94.}
Options, in order of preference: choose the metric and threshold properly (PR curves, cost-based threshold,
\Cref{ch:logistic}); class weights (\code{class_weight="balanced"}, \code{scale_pos_weight}); resampling ---
random under-sampling of the majority, random over-sampling, or \vocab{SMOTE} (synthetic minority points
interpolated between a minority point and one of its minority neighbours):
\codefile[code/ml/c15\_dataprep.py: SMOTE in ten lines]{gen/ml/snip/c15_dataprep-smote.py}
\begin{example}[SMOTE leakage]
	1000 rows, 5\% positives, 3-NN classifier. SMOTE the whole dataset, then cross-validate: F1
	$\gv{c15_dataprep/smote_wrong}$. SMOTE only inside each training fold, evaluate on untouched folds: F1
	$\gv{c15_dataprep/smote_right}$. In the leaky version synthetic points sit in the validation folds between
	their own parents in the training folds, and 3-NN finds them trivially.
\end{example}
\begin{remark}
	\trap Resample \emph{only the training data}, never the validation or test data (their class balance must
	match deployment), and remember resampling changes predicted probabilities. SMOTE also interpolates in
	feature space, which is meaningless for categorical features and risky in high dimensions.
\end{remark}

\section{Text features and TF-IDF}
\prototype{Three short documents; why does ``the'' get a low weight and ``cat'' a high one?}
Bag-of-words counts word occurrences (optionally $n$-grams). \vocab{TF-IDF} reweights a term's count in a
document by its rarity across the corpus:
\[ \text{tf-idf}(t,d)=\text{tf}(t,d)\cdot\IDF(t),\qquad \IDF(t)=\log\frac{N}{\text{df}(t)}\ \ \text{(textbook)}, \]
with $N$ documents and $\text{df}(t)$ the number containing $t$. scikit-learn's default is the smoothed
$\IDF(t)=\ln\frac{1+N}{1+\text{df}(t)}+1$, followed by L2-normalising each document vector.
\begin{example}[IDF by hand]
	Documents: ``the cat sat on the mat'', ``the dog sat on the log'', ``cats and dogs'' ($N=3$). ``the'' appears in
	$\gv{c15_dataprep/df_the}$ documents: textbook IDF $\ln(3/2)=\gv{c15_dataprep/idft_the}$, smoothed
	$\ln(4/3)+1=\gv{c15_dataprep/idfs_the}$. ``cat'' appears in $\gv{c15_dataprep/df_cat}$: $\ln3=
		\gv{c15_dataprep/idft_cat}$, smoothed $\ln2+1=\gv{c15_dataprep/idfs_cat}$ (note ``cats'' is a different token
	without stemming). In document 1, ``the'' occurs twice: unnormalised tf-idf $2\times1.2877=
		\gv{c15_dataprep/tfidf0_the}$ against $\gv{c15_dataprep/tfidf0_cat}$ for ``cat''. The smoothed IDFs and the
	normalised matrix match \code{TfidfVectorizer}. With $\log_{10}$ the textbook IDF of ``cat'' would be
	$\gv{c15_dataprep/idf_log10_cat}$: always state the base.
\end{example}
TF-IDF vectors are sparse and high-dimensional, compared with cosine similarity and fed to linear models or
naive Bayes; they ignore word order beyond $n$-grams and know no synonyms --- the gap embeddings fill.

\section{Time series in one page}
\prototype{Forecast tomorrow's orders from the last 30 days.}
A time series is not i.i.d.: values are autocorrelated, may trend and have seasonality. \vocab{Stationarity}
(constant mean, variance and autocorrelation over time) is assumed by classical models; differencing
$y_t-y_{t-1}$ removes trends, seasonal differencing $y_t-y_{t-s}$ removes seasonality, and the augmented
Dickey--Fuller test checks for a unit root. \vocab{ARIMA}$(p,d,q)$ = autoregression on $p$ lags + $d$
differences + a moving average of $q$ past errors; read $p$ from the partial autocorrelation and $q$ from the
autocorrelation plots. For ML models, build lag features, rolling statistics (computed only from the past),
calendar features and holiday flags, then validate with forward-chaining splits (\Cref{ch:eval}); remember tree
ensembles cannot extrapolate a trend, so model differences or detrend first.

\section{Recommenders in one page}
\prototype{Five users, four items, a sparse rating matrix: predict the blanks.}
\vocab{Content-based} recommenders use item features; \vocab{collaborative filtering} uses the rating matrix
itself: neighbourhood methods (item--item cosine similarity of rating columns) or \vocab{matrix factorisation}
$R\approx\U\V\T$ with $k$-dimensional user and item factors, fitted on observed entries only
\cite{koren2009}:
\[ \min_{\U,\V}\sum_{(u,i)\text{ observed}}(r_{ui}-\uu_u\T\vv_i)^2+\lambda\bigl(\norm{\U}_F^2+\norm{\V}_F^2\bigr). \]
\vocab{Alternating least squares} fixes $\V$ and solves a ridge regression for each user, then swaps. On a
$5\times4$ toy matrix ($k=2$, $\lambda=0.1$) each user update equals \code{Ridge} on that user's ratings, the
regularised objective decreases every sweep, the training RMSE reaches $\gv{c15_dataprep/als_rmse}$, and the model
predicts $\gv{c15_dataprep/als_pred_02}$ for user 1's unseen item 3. Practicalities: biases per user and item,
implicit feedback (clicks, not ratings), the cold-start problem for new users/items (fall back to content or
popularity), and evaluation by ranking metrics (precision@$k$, NDCG) on time-split data.

\begin{moral}
	Fit every transform inside the training fold; scale for distances and penalties, not for trees; treat
	missingness and rare categories as information to be estimated carefully, not noise to be papered over.
\end{moral}

\section\problemhead

\begin{problem}[IDF]
	\onechili
	A corpus has 1000 documents; ``refund'' appears in 10, ``the'' in 990. Compute textbook IDF (natural log) for
	both, and the tf-idf of ``refund'' in a document where it occurs 3 times.
	\begin{hint}
		$\ln(N/\text{df})$.
	\end{hint}
	\begin{sol}
		$\ln100=\gv{c15_problems/p15a_r}$, $\ln(1000/990)=\gv{c15_problems/p15a_t}$; tf-idf $3\times4.605=
			\gv{c15_problems/p15a_tf}$.
	\end{sol}
\end{problem}

\begin{problem}[Smoothed target encoding]
	\onechili
	Global conversion rate 5\%. Pincode A has 2 visits and 1 conversion; pincode B has 400 visits and 60. Compute
	smoothed encodings with $m=20$.
	\begin{hint}
		$(n\bar y+m\bar y_{\text{global}})/(n+m)$.
	\end{hint}
	\begin{sol}
		A: $(1+20\cdot0.05)/22=\gv{c15_problems/p15b_a}$ (raw mean $0.5$ shrunk hard). B: $(60+1)/420=
			\gv{c15_problems/p15b_b}$ (barely shrunk from $0.15$).
	\end{sol}
\end{problem}

\begin{problem}[Scaling, true or false]
	\onechili
	(a) Standardising features changes a random forest's predictions. (b) $k$-means results depend on feature units.
	(c) PCA should be run on standardised data when features have different units. (d) Gradient descent on linear
	regression converges faster after standardising. (e) The scaler should be fitted on train and test together for
	consistency.
	\begin{hint}
		Only one statement about trees; one about leakage.
	\end{hint}
	\begin{sol}
		(a) false; (b) true; (c) true; (d) true (better conditioning); (e) false --- fit on training data only and
		apply the same transform to test data.
	\end{sol}
\end{problem}
